\documentclass[11pt,a4paper]{article}
\usepackage[margin=2.5cm]{geometry}
\usepackage{hyperref}
\usepackage{listings}
\usepackage{xcolor}
\usepackage{amsmath}
\usepackage{parskip}

\lstdefinelanguage{mlgs}{
  keywords={let, fn, new, prot, cast, lambda, as, if, then, else},
  keywordstyle=\bfseries,
  comment=[l]{//},
  morestring=[b]",
}
\lstset{
  language=mlgs,
  basicstyle=\ttfamily\small,
  frame=single,
  breaklines=true,
  columns=fullflexible,
}

\title{\textbf{ML-GS: An Implementation of Gradual Security Typing with References}}
\author{Eren Yenigül}
\date{June 2026}

\begin{document}
\maketitle

\section{Scope and Motivation}

This project implements ML-GS, the calculus introduced by Fennell and Thiemann in
\emph{Gradual Security Typing with References} (CSF 2013). ML-GS extends a monomorphic
ML core language with references and higher-order functions with gradual information-flow
control (IFC). Security casts allow programmers to mix statically and dynamically checked
regions in the same program, bridging the gap between the precision of static type systems
and the flexibility of dynamic monitoring.

The motivation is practical: static IFC type systems such as FlowCaml are often too
conservative for real programs. For example, a function whose safety depends on a runtime
flag (e.g., \texttt{isPrivileged}) cannot be verified statically even when it is
semantically safe. By inserting casts, programmers can confine dynamic checks to those
regions and keep the rest statically verified.

The paper's key technical challenge is handling references under gradual typing. Unlike
pure-lambda-calculus settings, references allow flow-sensitivity attacks, require invariant
subtyping, and complicate the boundary between static and dynamic code. The paper addresses
this with a distinction between \emph{access types} and \emph{current types} on heap cells,
and with a no-sensitive-upgrade (NSU) policy for dynamic IFC.

\section{Approach and Main Results}

The implementation is a Scala interpreter that covers the full ML-GS calculus from the paper.
It consists of three main components.

\paragraph{Parser.} A combinator parser that reads ML-GS source programs. The surface syntax
extends the paper's formal notation with named functions, multi-argument currying, sequencing,
and explicit cast syntax. Security level annotations are written as \texttt{[L]}, \texttt{[H]},
and \texttt{[?]}. 

\paragraph{Type checker.} Implements the typing rule $pc;\Gamma;M \vdash e : t$ (Figure 2
of the paper), including subtyping ($\prec$), compatibility ($\sim$), and the typing rules for
casts (\textsc{T-Cast}), protection (\textsc{T-Protect}), guard casts (\textsc{T-GuardCast}),
function guard casts (\textsc{T-FunCast}), and pointer casts (\textsc{T-RefCast}). The type
checker enforces the program counter constraint on write effects and reports source-location
errors.

\paragraph{Interpreter.} Implements the small-step reduction relation $e/\mu \xrightarrow{PC} r$
from Figures 5, 6, and 9 of the paper. This includes:
\begin{itemize}
  \item Lambda application with protect wrapping (\textsc{R-App}, \textsc{R-Protect}).
  \item Heap allocation, dereference, and assignment with NSU enforcement
        (\textsc{R-New}, \textsc{R-Deref}, \textsc{R-Asgn}, \textsc{R-Asgn-NSU-Fail}).
  \item Cast reduction with propagation into sub-values (\textsc{R-Cast-Sub},
        \textsc{R-Cast-To-Dyn}, cast propagation of Figure 7).
  \item Guard casts that adjust program counter constraints through assignments and
        function bodies (\textsc{R-GuardCast-New}, \textsc{R-GuardCast-Asgn},
        \textsc{R-GuardCast-App}, \textsc{R-FunCast}, \textsc{R-RefCast}).
  \item Blame propagation: failures carry blame labels that identify the responsible cast.
\end{itemize}

A REPL is also included for interactive exploration. It has a history feature, keeping the previous bindings and the heap state after each evaluation with a press of the enter key. 

\subsection*{Example}

The following program from Section~II of the paper is supported. Casting
\texttt{sendToFacebook} to a dynamic type allows it to be used in the untyped
\texttt{addPrivileged} context. A downcast of the result to \texttt{ref~Report\textsuperscript{L}}
enforces the static guarantee at the boundary.

\begin{lstlisting}
fn sendToFacebook(r : ref[L]<int[L]>) -[L]-> { r := 0 };

let w = sendToFacebook as (ref[?]<int[?]> -[?]-> unit[?])[L];
// w is now dynamically checked
\end{lstlisting}

Running the NSU violation example from Section~II raises a blame exception at runtime,
correctly identifying the cast location.

\section{Evaluation}

The implementation is evaluated via a test suite of 20 unit tests (\texttt{ProgramTests.scala})
covering:

\begin{itemize}
  \item Basic values, let bindings, and function application.
  \item Heap allocation, dereference, and assignment.
  \item Security level propagation through pointer dereference.
  \item Cast acceptance: dynamic upcasts, H-to-? and L-to-? casts.
  \item Cast rejection: downcast of an H value to L raises a blame error at runtime.
  \item Static type errors: allocating an L reference under an H program counter.
  \item NSU violation: assigning to a public cell under a secret program counter.
  \item Higher-order functions and guard casts.
\end{itemize}

All 20 tests pass. The example programs in the \texttt{examples/} directory, including
\texttt{function\_guard\_cast.mlgs} and \texttt{pointer\_cast.mlgs}, also execute correctly.

\section{Limitations}

\begin{itemize}
  \item \textbf{Two-level lattice only.} The implementation uses the L/H lattice from the paper.
        The paper notes results generalize to arbitrary discrete lattices, but this is not implemented.
  \item \textbf{No polymorphism.} The calculus and implementation are monomorphic, as in the paper.
  \item \textbf{Termination insensitivity.} The system proves termination-insensitive
        non-interference (TINI). Timing and termination channels are not addressed.
  \item \textbf{Blame precision.} As discussed in Section VIII of the paper, blame labels
        can be imprecise because labels are joined unconditionally in \textsc{R-Deref}. The
        extended pointer syntax that would recover precision (also described in Section VIII)
        is not implemented.
  \item \textbf{No inference.} Cast placement is entirely manual. The paper lists automatic
        cast inference as future work.
  \item \textbf{No FE strategy.} The implementation uses NSU for dynamic IFC. The faceted
        execution (FE) strategy mentioned in the paper is not included.
\end{itemize}

\section{Experience}

The most interesting part of the project was working through the reference cast mechanism.
The distinction between the access type of a pointer and the current type of a heap cell is
subtle: a cast changes only the pointer's access type, not the stored value, and the
interpreter must insert a reconciling cast at each dereference. Getting this right required
reading the typing and reduction rules together several times, and the guard cast rules
(Figures 8 and 9) that propagate program counter constraints through sub-expressions were
the most tricky to implement correctly.

The main challenge was the sheer number of mutually dependent concepts: security levels,
type annotations, blame labels, the PC context, and the heap current/access type split all
interact in almost every reduction rule. Building the interpreter rule by rule and validating
each with a targeted test case kept the process manageable. Overall the paper is well-structured,
and having the formal rules as a direct blueprint made the Scala code relatively straightforward
once the data model was settled.

\end{document}
